% uid: ""
% school: "lycee"
% level: "3eme"
% section: ["Economie et Gestion"]
% chapter: "Statistiques"
% title: "Augmentation de salaires : effet d'une transformation affine"
% difficulty: "Facile"
% points: null
% tags: ["statistiques", "transformation-affine", "moyenne", "ecart-type", "coefficient-de-variation"]
% provenance: "own"
% dressing: "transformation affine"
% rung: 1
% batch: "2026-11-20-stat-3eco-A"

% ---STATEMENT---
Dans une entreprise de $40$ employés, le salaire mensuel $x$ (en dinars) a pour moyenne $\overline{x}=600$ et pour écart type $\sigma_x=80$.

La direction accorde à chacun une augmentation de $5\,\%$, puis ajoute une prime fixe de $20$ dinars. Le nouveau salaire d'un employé qui gagnait $x$ dinars est noté $y$.

\begin{enumerate}
  \item Exprimer $y$ en fonction de $x$.
  \item Calculer la moyenne $\overline{y}$ et l'écart type $\sigma_y$ des nouveaux salaires.
  \item Calculer la variance $V_y$.
  \item Calculer les coefficients de variation $\dfrac{\sigma_x}{\overline{x}}$ et $\dfrac{\sigma_y}{\overline{y}}$. L'écart relatif entre les salaires a-t-il augmenté ou diminué ?
\end{enumerate}

% ---SOLUTION---
\begin{enumerate}
  \item Une augmentation de $5\,\%$ multiplie par $1+\dfrac{5}{100}=1{,}05$. Si le salaire est $x$, il devient $1{,}05\,x$, puis on ajoute la prime de $20$ dinars :
  $$ \ib{$y=1{,}05\,x+20$}$$
  \item \textbf{Rappel.}\par
  \begin{prop}[Transformation affine d'une série]
  Si $y=ax+b$ pour chaque valeur, alors
  $$\overline{y}=a\,\overline{x}+b\qquad\text{et}\qquad \sigma_y=|a|\,\sigma_x .$$
  La constante $b$ décale toutes les valeurs de la même quantité : elle ne change donc \emph{pas} la dispersion.
  \end{prop}
  Ici $a=1{,}05$ et $b=20$ :
  $$\overline{y}=1{,}05\times600+20=630+20=650 \qquad \sigma_y=1{,}05\times80=84.$$
  Le salaire moyen est maintenant de $650$ dinars et l'écart type de $84$ dinars.
  \item La variance est le carré de l'écart type :
  $$V_y=\sigma_y^2=84^2=7\,056.$$
  \emph{Vérification :} $V_y=a^2V_x=1{,}05^2\times80^2=1{,}1025\times6\,400=7\,056$ $\checkmark$
  \item Le coefficient de variation mesure la dispersion \emph{relative} : $CV=\dfrac{\sigma}{\overline{x}}$.
  $$\frac{\sigma_x}{\overline{x}}=\frac{80}{600}=\frac{2}{15}\approx0{,}133\quad(13{,}3\,\%)\qquad
  \frac{\sigma_y}{\overline{y}}=\frac{84}{650}\approx0{,}129\quad(12{,}9\,\%).$$
  La dispersion en dinars augmente ($80\to84$), mais la dispersion relative diminue légèrement ($13{,}3\,\%\to12{,}9\,\%$).

  \textbf{Pourquoi ?} Tous les salaires sont multipliés par $1{,}05$ (le $CV$ resterait inchangé), mais la prime fixe de $20$ dinars augmente la moyenne sans augmenter l'écart type : le $CV$ baisse.
\end{enumerate}

\textbf{Idée :} une transformation affine $y=ax+b$ déplace la moyenne comme elle déplace les valeurs, mais la constante $b$ ne change pas la dispersion.
